\subsection{Movilidad y capacidad de superación de terreno del prototipo}

\subsubsection{Descripción de la prueba}

Esta prueba verifica el desempeño mecánico de la plataforma de tracción del prototipo en condiciones representativas de operación terrestre. Se evalúan tres capacidades principales: desplazamiento en terreno plano, giro sobre eje vertical y superación de irregularidades del terreno.

\paragraph{Objetivos mínimos}
\begin{itemize}
    \item Sostener una velocidad de traslación de $2\ \text{m/s}$ durante $30\ \text{s}$ en terreno plano.
    \item Realizar una rotación sobre su eje de $180^\circ$ en $2\ \text{s}$.
    \item Trepar una pendiente de $15^\circ$ sin pérdida de estabilidad.
    \item Superar obstáculos de altura entre $5$ y $7\ \text{cm}$ y ancho de $10\ \text{cm}$.
\end{itemize}

\subsubsection{Criterios de aprobación}

\begin{table}[H]
\centering
\small
\begin{tabular}{p{8.8cm} p{4.1cm}}
\toprule
\textbf{Métrica} & \textbf{Criterio de aprobación} \\
\midrule
Velocidad media en tramo plano & $\geq 2{,}0\ \text{m/s}$ \\
Tiempo continuo de operación a velocidad objetivo & $\geq 30\ \text{s}$ \\
Rotación sobre eje propio & $180^\circ$ en $\leq 2{,}0\ \text{s}$ \\
Error final de orientación tras giro de $180^\circ$ & $\leq 5^\circ$ \\
Superación de pendiente & $15^\circ$ sin detención ni vuelco \\
Superación de obstáculo rectangular (5--7 cm alto, 10 cm ancho) & 3/3 intentos exitosos \\
\bottomrule
\end{tabular}
\caption{Criterios de aprobación de la prueba de movilidad del prototipo.}
\label{tab:criterios_prueba_movilidad}
\end{table}

\subsubsection{Procedimiento de ejecución}

\begin{enumerate}
    \item Preparar una pista de ensayo plana con longitud suficiente para sostener el régimen de velocidad durante 30 s.
    \item Instrumentar medición de velocidad y trayectoria (odometría, registro temporal y referencia externa).
    \item Ejecutar corrida en terreno plano y registrar velocidad instantánea y promedio.
    \item Ejecutar maniobra de rotación sobre eje propio de $180^\circ$ y medir tiempo de maniobra.
    \item Repetir la rotación al menos tres veces para verificar repetibilidad.
    \item Ejecutar ascenso por rampa calibrada a $15^\circ$.
    \item Ejecutar cruce de obstáculo rectangular (alto 5--7 cm, ancho 10 cm), con tres repeticiones.
    \item Consolidar resultados y verificar cumplimiento de criterios.
\end{enumerate}

\subsubsection{Resultados}

\begin{table}[H]
\centering
\small
\begin{tabular}{p{8.8cm} p{2.2cm} p{2.8cm}}
\toprule
\textbf{Métrica} & \textbf{Resultado} & \textbf{Cumple} \\
\midrule
Velocidad media en tramo plano [m/s] & --- & --- \\
Tiempo sostenido a velocidad objetivo [s] & --- & --- \\
Tiempo de giro 180$^\circ$ [s] & --- & --- \\
Error final de orientación [deg] & --- & --- \\
Superación de pendiente 15$^\circ$ [Sí/No] & --- & --- \\
Éxitos en cruce de obstáculo [\#/\#] & --- & --- \\
\bottomrule
\end{tabular}
\caption{Resultados medidos de la prueba de movilidad del prototipo.}
\label{tab:resultados_prueba_movilidad}
\end{table}
